\documentclass[10pt,a4paper]{amsart}
\usepackage[margin=19mm]{geometry}
\usepackage{amsmath,amssymb,mathtools}
\usepackage[scr=rsfs,cal=euler]{mathalfa}
\usepackage{xfrac}
\usepackage[unicode,hidelinks,bookmarks=false]{hyperref}
\newcommand{\ie}{i.e.\ }
\newcommand{\cf}{cf.\ }
\newcommand{\resp}{respectively\ }
\newcommand{\Subsection}{\subsection}
\providecommand{\nobreakdash}{\nobreak\hbox{-}}
\setlength{\emergencystretch}{2em}
\multlinegap0pt
\usepackage{siunitx}
\usepackage{siunitx}
\usepackage{siunitx}
\usepackage{siunitx}
\usepackage[shortlabels]{enumitem}
\usepackage{amssymb}
\usepackage{mathtools}
\usepackage{bm}
\usepackage{siunitx}
\usepackage{xcolor}
\usepackage{tikz}
\newtheorem{proposition}{Proposition}
\newcommand{\Jop}{\bm J}
\newcommand{\Kop}{\bm K}
\newcommand{\R}{\mathbb{R}}
\newcommand{\Rop}{\bm R}
\title{An Energy-Based Conservative--Dissipative Latent Neural Evolution Operator for Magnetization Dynamics}
\author{Sebastian Schaffer, Lukas Exl}
\date{Source: arXiv:2609.04530v1 (CC BY 4.0)}
\begin{document}
\maketitle

\noindent\emph{Source and licence.} An Energy-Based Conservative--Dissipative Latent Neural Evolution Operator for Magnetization Dynamics. Version arXiv:2609.04530v1 (CC BY 4.0), distributed by arXiv under the Creative Commons Attribution 4.0 International licence, http://creativecommons.org/licenses/by/4.0/, as recorded in the arXiv licence metadata for this exact version and shown on the abstract page https://arxiv.org/abs/2609.04530v1. Full licence notice: \texttt{licenses/arxiv-2609.04530-CC-BY-4.0.txt}.\par\smallskip

\noindent\textit{Editorial scope.} Counted unit: the article's proposition prop:operator\_spectra (source0.tex lines 325--336). The article's complete original proof of it is delivered with it (source0.tex lines 338--356, one \texttt{proof} environment of 19 lines). No mathematics is written, repaired or re-derived: the statement and the proof are the article's own lines, in the article's own notation, and no retained line is substituted or re-flowed.  No further source-local context is needed: every cross-reference of every retained line resolves inside the two delivered spans.  Nothing was omitted from any retained span, so the package carries no omission record.  No retained line contains a citation, so the wrapper carries no bibliography and no bibliography entry.  No retained line contains a figure environment, an image-inclusion command or drawing code, so no image is delivered and the package declares no asset dependency.  Editorial formatting only, in addition: an AMS \texttt{amsart} preamble that restates the article's own declarations and macros used by the retained lines, namely the environments \texttt{proposition} on separate counters, together with the standard packages those lines need.  The retained environments print in this document as Proposition 1 in order of appearance, whereas the article prints the same items with the numbering of its own full document.  The complete native source of the article as distributed in the arXiv e-print for this exact version is bundled unchanged as \texttt{sources/209370/source0.tex}.
\begin{proposition}[Operator spectra]
Let $\Jop,\Rop\in\R^{d\times d}$ satisfy
$\Jop^{\mathsf T}=-\Jop$ and $\Rop\succeq0$, and define
$\Kop=\Jop-\Rop$. Then:
\begin{enumerate}[label=(\roman*)]
  \item every eigenvalue of $\Jop$ is purely imaginary or zero;
  \item every eigenvalue of $\Rop$ is real and nonnegative; and
  \item every eigenvalue $\lambda$ of $\Kop$ satisfies
        $\operatorname{Re}(\lambda)\leq0$.
\end{enumerate}
\label{prop:operator_spectra}
\end{proposition}

\begin{proof}
The first two statements are standard consequences of real antisymmetry and
real symmetric positive semidefiniteness. For the third, let
$\bm v\in\mathbb C^d\setminus\{\bm0\}$ satisfy
$\Kop\bm v=\lambda\bm v$, and let $\bm v^\ast$ denote its conjugate transpose.
Then
\begin{equation}
  \lambda
  =\frac{\bm v^\ast\Jop\bm v}{\bm v^\ast\bm v}
   -\frac{\bm v^\ast\Rop\bm v}{\bm v^\ast\bm v}.
\end{equation}
The first quotient is purely imaginary, while the second is real and
nonnegative. Therefore
\begin{equation}
  \operatorname{Re}(\lambda)
  =-\frac{\bm v^\ast\Rop\bm v}{\bm v^\ast\bm v}
  \leq0.
\end{equation}
\end{proof}

\end{document}
